\section{Phát biểu bài toán và các bất biến thiết kế}\label{sec:problem}

\subsection{Đối tượng và ký hiệu}
Một \emph{chương} được biểu diễn bằng đồ thị có hướng $G=(V,E)$. Mỗi đỉnh $v\in V$ là một \emph{tiết} học; mỗi cạnh $(u,v)\in E$ mang một trong hai kiểu thứ tự, \textsc{prerequisite} (học $u$ trước $v$) hoặc \textsc{application} ($v$ vận dụng $u$). Đỉnh $v$ đi kèm gói nội dung $K_v=(\mathcal{T}_v,\mu_v,\mathcal{W}_v,\mathcal{N}_v,\mathcal{I}_v,\mathcal{R}_v)$ gồm các cách giảng lý thuyết, phương pháp, ví dụ mẫu, ghi chú lỗi sai, ngân hàng câu hỏi và các ý chính của bài kể lại. Một \emph{câu hỏi} $i\in\mathcal{I}_v$ có đề, đặc tả đáp án $\alpha_i$ (số, trắc nghiệm, so sánh, sắp xếp, đúng/sai hoặc mở), mức nhận thức $\ell_i\in\{A,B,C,D,E,\mathrm{DIAG}\}$, độ khó $\delta_i\in[0,1]$, mục đích (kiểm tra hiểu, luyện tập, kiểm tra cuối tiết, ôn tập) và một thang gợi ý.

Một buổi học sinh ra \emph{sổ sự kiện} $L=(e_1,\dots,e_T)$ chỉ ghi thêm. Ở lượt $t$, học sinh nói $u_t$; hệ thống đọc $r_t=\mathrm{read}(u_t,s_t)$ (đáp án, xin gợi ý, hỏi lại, lạc đề, cảm xúc\dots), quyết định $d_t=\pi(s_t,r_t)$ gồm nước đi, bậc gợi ý, mục tiêu gợi ý và các mã luật, rồi diễn đạt $y_t=\varphi(d_t,K,P_t)$ với các đoạn nguồn $P_t$, và chỉ gửi $y_t$ khi bộ kiểm chứng $\mathcal{V}$ chấp nhận. Mức hiểu bài là $M=\Phi(L;\theta)$, một hàm thuần túy của sổ sự kiện và bộ tham số có phiên bản $\theta$. Có ba tầng quyết định: chính sách dạy trên đồ thị (luật \rid{N-*}, chọn tiết và hoạt động), chính sách một câu hỏi (luật \rid{R-*}, chọn nước đi và bậc gợi ý) và diễn đạt (mô hình ngôn ngữ). Chỉ tầng cuối không tất định.

\subsection{Sáu bất biến}
Từ các rủi ro ở Mục~\ref{sec:intro} chúng tôi rút ra sáu bất biến mà hệ thống phải giữ \emph{bất kể mô hình ngôn ngữ nói gì}; Bảng~\ref{tab:invariants} nêu cơ chế thực thi và cách kiểm chứng của từng bất biến. Hai nguyên tắc kỹ thuật bổ sung đi cùng: ràng buộc của buổi học được thực thi ở máy chủ chứ không bao giờ chỉ ẩn nút ở giao diện, và khách không bao giờ được chọn mô hình.

\begin{table}[tbp]
\centering
\caption{Sáu bất biến của CAVAT, cơ chế giữ chúng và nơi bài báo kiểm chứng.}
\label{tab:invariants}
\small
\begin{tabularx}{\linewidth}{@{}>{\raggedright\arraybackslash}p{2.7cm}YY@{}}
\toprule
Bất biến & Cơ chế thực thi & Kiểm chứng \\
\midrule
\textbf{B1.} Bí mật đáp án: đáp án câu kiểm tra không có trong ngữ cảnh gửi mô hình & Bộ dựng lời nhắc không đưa đáp án của câu kiểm tra; bậc gợi ý tối đa H1 trong kiểm tra (\rid{R-ASSESS-CAP}); chỉ nước đi \textsc{full\_explanation} được phép nói đáp án & Quét đáp án trong mọi lượt (Mục~\ref{sec:res-soc}, \ref{sec:market}) \\
\textbf{B2.} Chính sách quyết định: nước đi, bậc và mục tiêu gợi ý do luật chọn, mô hình chỉ diễn đạt & Hàm quyết định tất định có phiên bản; mô hình chỉ được gọi cho một tập nước đi; mọi quyết định ghi mã luật & Tần suất luật, số lượt do mô hình viết (Mục~\ref{sec:res-soc}) \\
\textbf{B3.} Chấm tất định: đúng/sai Toán do bộ chấm, không do mô hình & Bộ chấm viết lại theo chuẩn vi-VN: dấu phẩy, số đọc bằng chữ, dạng đề hỏi, ý bắt buộc của đáp án mở & Bộ chấm tay giữ riêng (Mục~\ref{sec:market}) \\
\textbf{B4.} Điểm tái lập: điểm là phép chiếu từ sổ sự kiện; mô hình không ghi điểm & Hàm gấp sự kiện thuần túy; nhãn do mô hình chỉ là bằng chứng có trọng số & Kiểm tương đương và phát lại (Mục~\ref{sec:res-assess}) \\
\textbf{B5.} Phạm vi: không dạy trước nội dung của tiết sau tiết đang học; câu ngoài chương được giữ lại cho giáo viên & Phạm vi truy xuất chỉ gồm tiết đang học và tổ tiên trên đồ thị; so khớp tiết sau trước khi trả lời & Câu hỏi bài sau, so với RAG phẳng (Mục~\ref{sec:res-rag}) \\
\textbf{B6.} Nền tảng: mọi điều gia sư khẳng định có trong đoạn nguồn hoặc lời học sinh & Kiểm chứng số, thuật ngữ Toán, từ có nghĩa; thử lại một lần rồi dùng câu trích nguyên văn hoặc giữ cho giáo viên & Độ bám nguồn, số lần bịa (Mục~\ref{sec:res-rag}) \\
\bottomrule
\end{tabularx}
\end{table}

\subsection{Ràng buộc của môi trường}
CAVAT được thiết kế để chạy trên các dịch vụ miễn phí, nên giới hạn hạn mức là một ràng buộc thiết kế chứ không phải chi tiết triển khai. Mỗi lời nhắc gửi mô hình tutor không quá $3\,000$ token và tắt suy luận mở rộng, vì hạn mức $8\,000$ token mỗi phút; mỗi buổi học có trần $15$ lần gọi mô hình; khi mô hình không dùng được (hết hạn mức, mạng, lỗi), mọi nước đi vẫn có câu mẫu đã duyệt để nói thay, nên buổi học không bao giờ dừng vì nhà cung cấp. Hệ quả là một phần đáng kể lời gia sư trong các buổi dài là câu mẫu của giáo viên; Mục~\ref{sec:res-soc} đo tỉ lệ này và Mục~\ref{sec:limits} nêu ý nghĩa của nó cho cách đọc kết quả.
